\section{Background}

\label{section:dynamical-geometric-background}

Throughout this section, we~let $(\Sigma, g)$ denote a closed, connected, oriented hyperbolic surface—that is, a Riemannian surface with constant curvature $-1$. We~establish the following results:

\begin{itemize}
\item In Section~\ref{ssection:geo}, we~recall several technical results concerning the geometry and dynamics on the tangent bundle $T\Sigma$. In~particular, we~emphasize the unique ergodicity of the horocyclic flow, a key ingredient in the proof of Theorem~\ref{theorem:horocyclic}.

\item In Section~\ref{ssection:magnetic-flow}, we~introduce the magnetic flows on both the tangent and cotangent bundles, and show that they are equivalent under the identification induced by the metric. We~also prove a technical result on the propagation of functions along the magnetic flow, which will be instrumental in the proof of Theorem~\ref{theorem:horocyclic}.

\item Finally, in Section~\ref{ssection:landau}, we~review standard properties of magnetic Laplacians. In~particular, we~explicitly compute the bottom of their spectrum—the so-called Landau levels. \end{itemize}

\subsection{Geometry and dynamics on the tangent bundle}

\label{ssection:geo}
Let $T\Sigma$ be the tangent bundle of $\Sigma$ and
\begin{equation}
\label{equation:footpoint}
\pi : T\Sigma \to \Sigma
\end{equation}
be the footpoint projection. We~refer the reader to \cite[Ch.\,1]{Paternain-99} or \cite[Ch.\,13]{Lefeuvre-book} for background material on the geodesic flow.

\subsubsection{Structural equations} The geodesic flow $\varphi_t : T\Sigma \to T\Sigma$ is defined as
\[
\varphi_t(x,v) := (\gamma(t), \dot{\gamma}(t)),
\]
where $\gamma : \R \to \Sigma$ is the (unique) curve solving Newton's equation:
\begin{gather} \label{eq:newton_geodesic}
\nabla_{\dot{\gamma}(t)}\dot{\gamma}(t) = 0, \qquad \gamma(0)=x, \; \dot{\gamma}(0)=v,
\end{gather}
and $\nabla$ is the Levi-Civita connection. We~let
\[
X := \partial_t\varphi_t|_{t=0} \in C^\infty(T\Sigma, T(T\Sigma))
\]
be its generator.

As $\Sigma$ is oriented, there is a well-defined fiberwise rotation $R_\theta : T\Sigma \to T\Sigma$ by angle $\theta \in [0,2\pi)$ in the fibers of $T\Sigma$. Let $V$ be the generator of this rotation, namely
\[
V := \partial_\theta R_\theta|_{\theta = 0}.
\]
Let $V_\perp$ be the Euler vector field of $T\Sigma$, i.e. the generator of the flow
\begin{equation}
\label{equation:flot-vperp}
\varphi_t^{V_\perp}(x,v) = (x,e^tv).
\end{equation}
Finally, define
\[
X_\perp := [V,X].
\]
The vector fields $\{X,X_\perp,V,V_\perp\}$ form a basis of $T(T\Sigma)$ at any point of $T\Sigma$. The metric for which this is an orthonormal frame is called the \emph{Sasaki metric}.

We will need the following result:

\begin{lemma}
The above vector fields satisfy the commutation relations:
\begin{equation}
\label{equation:commutator-relations}
\begin{aligned}
[X,X_\perp] &= - |v|^2 V,& [X,V] &=-X_\perp, & [X_\perp,V] &= X, \\
[V_\perp,X] &= X, & [V_\perp, X_\perp] &= X_\perp, & [V,V_\perp] &= 0.
\end{aligned}
\end{equation}
\end{lemma}

We refer to \cite[Lem.\,15.2.1]{Lefeuvre-book} for a proof. (In this reference, $X_\perp$ is denoted by~$H$; the computations are done on the unit tangent bundle but the above relations follow easily by a scaling argument.) Notice that in our case, the sectional curvature is $\kappa = -1$.

\subsubsection{Horocyclic flow} For $\lambda \geq 0$, define
\[
S^\lambda\Sigma := \{|v|_g=\lambda\}.
\]
We let $S\Sigma := S^{1}\Sigma$ be the unit tangent bundle. It is straightforward to verify that $X,X_\perp$ and $V$ preserve the layers $S^\lambda\Sigma$.

The geodesic flow $(\varphi_t)_{t \in \R}$ is Anosov on $S\Sigma$ in the sense that {\color{black}there exists a continuous flow-invariant decomposition}
\[
T(S \Sigma) = \R X \oplus E_s \oplus E_u,
\]
where
\[
\begin{split}
|d\varphi_t(w)| \leq e^{-t} |w|, &\qquad \forall t \geq 0,\; w \in E_s \\
|d\varphi_{-t}(w)| \leq e^{-t} |w|, &\qquad \forall t \geq 0,\; w \in E_u.
\end{split}
\]
Here $|\csbullet|$ denotes the norm induced by the \emph{Sasaki metric} on $T\Sigma$, which is defined as the metric for which $\{X,X_\perp,V,V_\perp\}$ is an orthonormal frame.

We define the (stable) horocyclic vector field by
\[
U_+ := X_\perp - V.
\]
Let $(h_t)_{t \in \R}$ be the (stable) \emph{horocyclic flow} generated by $U_+$ on $S\Sigma$. Recall that the Liouville measure on $S\Sigma$ is the unique volume form $\mu_{\mathrm{Liouv}}$ such that $\mu_{\mathrm{Liouv}}(X,X_\perp,V) = \mathrm{cst}$. The constant is normalized such that $\mu_{\mathrm{Liouv}}$ is a probability measure on $S\Sigma$. The Liouville measure is invariant by $X, X_\perp$ and $V$; it is thus also invariant by $U_+$.

\begin{theorem}
\label{lemma:horocycle}
The following holds:
\begin{enumerate}
\item\label{lemma:horocyclei} The flow $(h_t)_{t \in \R}$ is uniquely ergodic on $S\Sigma$ and $\mu_{\mathrm{Liouv}}$ is the only flow-invariant probability measure.

\item\label{lemma:horocycleii} Let $\theta\in(0,1/2)$ be such that $\theta(1-\theta) \leq \lambda_1(\Sigma)$, where $\lambda_1(\Sigma) > 0$ is the first non-zero eigenvalue of the Laplacian $\Delta_g$ on functions. Then there exists $C > 0$ such that for all $a \in C^{\infty}(S\Sigma)$, for all $T > 0$,
\[
\sup_{v \in S\Sigma}\, \biggl|\dfrac{1}{T}\int_0^T a(h_t(v))~ \dd t - \int_{S\Sigma} a(v) ~\dd\mu_{\mathrm{Liouv}}\biggr| \leq C T^{-\theta} \|a\|_{H^3(S\Sigma)},
\]
where $H^3(S\Sigma)$ denotes the $L^2$-based Sobolev norm of order $3$.
\end{enumerate}
\end{theorem}

Unique ergodicity in constant curvature was established by Furstenberg \cite{Furstenberg-73}, while the rate of convergence of ergodic averages was proved by Burger \cite{Burger-90}. A~refinement of this result can be found in \cite{Flaminio-Forni-03}.

\subsection{Magnetic hyperbolic flow}

\label{ssection:magnetic-flow}

In this paragraph, we~discuss the magnetic flow both on the tangent and the cotangent bundles.

\subsubsection{Magnetic flow on the tangent bundle}
\label{sssection:magnetic-tsigma}
A {\em magnetic geodesic}
is a curve $\ga $ of $\Sigma$ which satisfies the equation
\begin{equation} \label{eq:Newton_Lorentz}
\nabla_{\dot{\gamma}(t)} \dot{\gamma}(t) = - B j_{\gamma(t)}\dot{\ga} (t),
\qquad \gamma(0)=x,\; \dot{\gamma}(0)=v,
\end{equation}
where $(x,v) \in T\Sigma$, $\nabla$ stands for the Levi-Civita covariant derivative, $j$ is the almost-complex structure and
$ B \in C^\infty(\Sigma)$ is the magnetic {\em intensity}. {\color{black}As $\Sigma$ is oriented, and equipped with a metric $g$, the almost-complex structure $j \in C^\infty(\Sigma,\mathrm{End}(T\Sigma))$ is the fiberwise endomorphism of $T\Sigma$ given by rotation of tangent vectors by an angle $+\pi/2$.} In the sequel we assume that $B$ is constant
and positive.
The {\em magnetic flow} $(\Phi_t)_{t \in \R}$ is the autonomous flow of $T \Si$ such that for
any curve $\ga$ satisfying \eqref{eq:Newton_Lorentz},
\[
\Phi_t (x,v) := (\ga (t), \dot {\ga} (t)), \qquad \forall t \in \R.
\]
\begin{lemma} The magnetic flow is generated by the vector field $F := X - B V$.
\end{lemma}

\begin{proof} We work in a chart $U$ of $\Si$ that we identify with an open
subset of $\R^2$, so~$TU = U \times \R^2$. Let $(x,v) \in TU$. Let $c_1$,
$c_2$ be a geodesic and a magnetic geodesic respectively such that $ c_1
(0) = x = c_2 (0)$, $\dot{c}_1 (0) =v = \dot{c}_2 (0)$. Then $X(x,v) = (
\dot{c}_1 (0), \ddot{c}_1 (0))$ and $F(x,v) = (
\dot{c}_2 (0), \ddot{c}_2 (0)) $, thus
\[
F(x,v) - X(x,v) = (0, \ddot{c}_2 (0) - \ddot{c}_1 (0)) = (0,
(\nabla_{\dot{c}_2} \dot{c}_2) (0) - (\nabla_{\dot{c}_1} \dot{c}_1) (0)
) = (0, -B j_{x} v),
\]
where we have used that $\nabla_{\dot{c}_i} \dot{c}_i = \ddot{c}_i +
\Gamma_{c_i} (\dot{c}_i) (\dot{c}_i)$, $\Gamma$ being the connection
one-form, and then Newton's equations \eqref{eq:newton_geodesic},
\eqref{eq:Newton_Lorentz}. To conclude use that $V (x,v) = (0, j_x v)$.
\end{proof}

\begin{remark}The flow $(\Phi_t)_{t \in \R}$ is considered here on the tangent bundle $T\Sigma$, whereas it was initially defined on the cotangent bundle $T^*\Sigma$ in the introduction. In~Section~\ref{sssection:flow-cotangent}, we~show that these two flows are equivalent under the musical isomorphism induced by the metric. To avoid unnecessary notation, unless specified explicitly, we~do not distinguish between the flows on $T\Sigma$ and $T^*\Sigma$ in what follows.
\end{remark}

Recall that $(\varphi_t)_{t \in \R}$ is the geodesic flow (generated by $X$), $(R_t)_{t \in \R}$ is the $2\pi$-periodic rotation in the circle fibers of $S\Sigma$ (generated by $V$).

\begin{proposition}
\label{proposition:dynamic}
For any $\lambda> 0$, the following holds:
\begin{enumerate}

\item\label{proposition:dynamici} If $ \lambda< B$, $(\Phi_t)_{t \in \R}$ on $S^\lambda\Sigma$ is conjugate to $(R_{t/T_\lambda})_{t \in \R}$ on $S\Sigma$, where $T_\lambda = (B^2 - \lambda^2)^{-\sfrac{1}{2}}$;

\item\label{proposition:dynamicii} If $ \lambda = B$, then $(\Phi_t)_{t \in \R}$ on $S^\lambda\Sigma$ is conjugate to $(h_t)_{t \in \R}$ on $S\Sigma$;

\item\label{proposition:dynamiciii} If $ \lambda>B$, then $(\Phi_t)_{t \in \R}$ on $S^\lambda\Sigma$ is conjugate to $(\varphi_{t/T'_\lambda})_{t \in \R}$ on $S\Sigma$, where $ T_\lambda' = (\lambda^2 -B^2)^{-\sfrac{1}{2}}$.
\end{enumerate}
\end{proposition}

\begin{proof}
For $\lambda > 0$, let $\Psi_\lambda : T\Sigma \to T\Sigma$ be the map defined by $\Psi_\lambda(x,v) := (x,\lambda v)$. Observe that $\Psi_\lambda^* X = \lambda X$ and $\Psi_\lambda^*V=V$. The flow of $F = X-BV$ on $S^\lambda \Sigma$ is thus conjugate to the flow of $\Psi_\lambda^* F = \lambda X -BV$ on $S\Sigma$.

{\color{black}As $(\Sigma,g)$ is a hyperbolic surface, there exists a discrete torsion-free subgroup $\Gamma < \mathrm{PSL}(2,\R)$ such that $S\Sigma = \Gamma \backslash S\HH^2 = \Gamma\backslash\mathrm{PSL}(2,\R)$.} The vector field $\lambda X - BV$ is represented by the following matrix element of $\mathfrak{sl}(2,\R)$:
\[
\arraycolsep3.5pt
\lambda \begin{pmatrix} 1/2 & 0 \\ 0 & -1/2 \end{pmatrix} - B \begin{pmatrix} 0 & 1/2 \\ -1/2 & 0\end{pmatrix} = \begin{pmatrix} \lambda/2 & -B/2 \\ B/2 & -\lambda/2 \end{pmatrix}.
\]
For $\lambda = 1, B=0$, this is the generator of the geodesic flow; for $\lambda=0$, $B=1$, this is the generator of the rotation flow in the circle fibers, see \cite[\S 2.1]{Flaminio-Forni-03} for details. The eigenvalues of this matrix are $\pm\tfrac{1}{2}(\lambda^2-B^2)^{1/2}$ if $\lambda \geq B$ and $\pm \tfrac{i}{2}(B^2-\lambda^2)^{1/2}$ if $B \geq \lambda$. Notice that for $\lambda = B$, one finds the matrix
\[
\arraycolsep3.5pt
\lambda/2 \begin{pmatrix} 1 & -1 \\ 1 & -1 \end{pmatrix}
\]
which is conjugate by an element of $\mathrm{PSL}(2,\R)$ to
\[
\arraycolsep3.5pt
\begin{pmatrix} 0 & 1 \\ 0 & 0 \end{pmatrix},
\]
the generator of the horocyclic flow.

If two matrices in $\mathfrak{sl}(2,\R)$ are conjugate (by an element of $\mathrm{PSL}(2,\R)$), then their corresponding flows are conjugate too. Hence, if~$B > \lambda$, $(\Phi_t)_{t \in \R}$ on $S^\lambda\Sigma$ is conjugate to $(R_{t(B^2-\lambda^2)^{1/2}})_{t \in \R} = (R_{t/T_\lambda})_{t \in \R}$ on $S\Sigma$. If $B=\lambda$, $(\Phi_t)_{t \in \R}$ on $S^\lambda\Sigma$ is conjugate to $(h_t)_{t \in \R}$ on $S\Sigma$. If $B < \lambda$, $(\Phi_t)_{t \in \R}$ on $S^\lambda\Sigma$ is conjugate to $(\varphi_{t(\lambda^2-B^2)^{1/2}})_{t \in \R} = (\varphi_{t/T'_\lambda})_{t \in \R}$ on $S\Sigma$.
\end{proof}

\subsubsection{Magnetic flow on the cotangent bundle} \label{sssection:flow-cotangent}

We now define the magnetic flow on the cotangent bundle, and show that it coincides with that on the tangent bundle up to the identification provided by the metric. Let
\[
\flat : T\Sigma \to T^*\Sigma, \qquad \flat(x,v) := (x, g_x(v,\csbullet))
\]
be the musical isomorphism. Let $A$ be the Liouville
form on $T^*\Sigma$ defined by
\begin{equation}
\label{equation:liouville}
A_{(x, \xi) } := \xi (\dd_{x, \xi} \pi (\csbullet)),
\end{equation}
where $\pi : T^*\Sigma \to \Sigma$ denotes the projection. Consider the symplectic form
\[
\Om \in C^\infty(T^* \Si,\Lambda^2 T^*(T^*\Sigma)), \qquad \Om := dA + B~ \pi^*\op{vol},
\]
where $\op{\vol} \in C^\infty(\Si,\Lambda^2 T^*\Sigma)$ the Riemannian volume. Let $p(x,\xi) := \tfrac{1}{2} | \xi|^2$, and define the flow $(\Phi_t^{T^*\Sigma})_{t \in \R}$ as the Hamiltonian flow of the function $p \in C^\infty(T^*\Sigma)$ with respect to the symplectic form $\Omega$. Namely $(\Phi_t^{T^*\Sigma})_{t \in \R}$ is generated by $H^{\Omega_p}$ such that
\[
dp = \Omega(\csbullet, H^\Omega_p).
\]
To avoid confusion, let us denote temporarily $(\Phi_t^{T\Sigma})_{t \in \R}$ the magnetic flow defined on $T\Sigma$ in Section~\ref{sssection:magnetic-tsigma}.

\begin{lemma}
The map $\flat$ intertwines $(\Phi_t^{T\Sigma})_{t \in \R}$ and $(\Phi_t^{T^*\Sigma})_{t \in \R}$, that is
\[
\Phi_t^{T^*\Sigma} \circ \flat = \flat \circ \Phi_t^{T\Sigma}, \qquad \forall t \in \R.
\]
\end{lemma}

\begin{proof}
Define the $1$-form $\alpha := \flat^*A$ on $T\Sigma$. It satisfies
\[
\alpha_{(x,v)} = g_x(\dd_{x,v}\pi(\csbullet),v),
\]
see \cite[Lem.\,1.37, item\,1]{Paternain-99}. Hence
\[
\flat^*\Omega = d\alpha + B~ \pi^*\vol,
\]
where $\pi : T\Sigma \to \Sigma$ is the footpoint projection. Define $h(x,v) := |v|^2_g/2$. Let $H_h^{\flat^*\Omega}$ be the Hamiltonian vector field of $h$, computed with respect to $\flat^*\Omega$ on $T\Sigma$, namely
\[
dh = \flat^*\Omega(\csbullet, H_h^{\flat^*\Omega}).
\]
It is immediate that $\flat^* H_p^{\Omega} = H_h^{\flat^*\Omega}$ so the claim boils down to showing that
\[
H_h^{\flat^*\Omega} = X - B V.
\]
For that, we~compute:
\begin{equation}
\label{equation:mouch}
\flat^*\Omega(\csbullet,X-BV) = \dd\alpha(\csbullet, X) - B \dd \alpha(\csbullet, V) + B \pi^*\vol(\csbullet, X).
\end{equation}
A quick computation reveals that $\dd \alpha(\csbullet, X) = dh$. Indeed, $\dd \alpha(Y,X) = 0 = dh(Y)$ for $Y = X,X_\perp,V$ and $\dd \alpha(V_\perp,X) = |v|^2 = dh(V_\perp)$ using the $2$-homogeneity of $h$ and that $V_\perp$ is the Euler vector field.

In addition, using \cite[Prop.\,1.24]{Paternain-99} for the formula for $\dd \alpha$, we~find:
\[
\dd \alpha(\csbullet,V) = - g_x(j(x)v, \dd \pi(\csbullet)) = \pi^*\vol(\csbullet,X),
\]
so the last two terms in \eqref{equation:mouch} cancel out. This proves that $\flat^*\Omega(\csbullet,X-BV) = dh$, and thus $H_h^{\flat^*\Omega} = X - B V$.
\end{proof}
